% fig_gauge_innerproduct.tex  --  canonical forms and pairings.
% Requires: \usepackage{tikz}; \usetikzlibrary{arrows.meta, positioning, calc, fit, backgrounds}.
\begin{figure*}[t]
    \centering
    \begin{tikzpicture}[
            font=\scriptsize,
            >={Stealth[length=3pt]},
            % TN conventions as in Figs. 2-3: bonds 0.5pt, physical legs 0.7pt
            bond/.style ={line width=0.5pt},
            phys/.style ={line width=0.7pt},
            tnA/.style  ={draw, fill=black, circle, inner sep=1.6pt},
            tnAc/.style ={draw, fill=white, circle, inner sep=1.6pt},
            opf/.style  ={draw, fill=black!15, minimum size=0.34cm, inner sep=0pt},
            opc/.style  ={draw, fill=white, minimum size=0.34cm, inner sep=0pt},
            nlabel/.style={font=\scriptsize\bfseries, align=center},
            note/.style ={font=\scriptsize\itshape, text=black!70, align=center},
            x=1cm, y=1cm
        ]

        %% ===== Panel (a): GAUGE (TN) =====
        \node[nlabel, text width=3.6cm] at (1.4,5.4) {(a) gauge of $E_A$};

        \begin{scope}[shift={(0,4.0)}]
            \node[tnA] (A) at (0,0.28) {}; \node[tnAc] (Ac) at (0,-0.28) {};
            \draw[phys] (A)--(Ac);
            \draw[bond] (A.west)--(-0.55,0.28); \draw[bond] (Ac.west)--(-0.55,-0.28);
            \draw[bond] (A.east)--(0.4,0.28);   \draw[bond] (Ac.east)--(0.4,-0.28);
            \draw[bond] (0.4,0.28) to[out=0,in=0] (0.4,-0.28);
            \node at (1.15,0) {$=$};
            \draw[bond] (1.5,0.28)--(2.0,0.28); \draw[bond] (1.5,-0.28)--(2.0,-0.28);
            \draw[bond] (2.0,0.28) to[out=0,in=0] (2.0,-0.28);
        \end{scope}
        \node[nlabel] at (1.1,3.35) {right-canonical (unital)};

        \begin{scope}[shift={(0,2.0)}]
            \node[tnA] (B) at (0,0.28) {}; \node[tnAc] (Bc) at (0,-0.28) {};
            \draw[phys] (B)--(Bc);
            \draw[bond] (B.east)--(0.55,0.28); \draw[bond] (Bc.east)--(0.55,-0.28);
            \draw[bond] (B.west)--(-0.4,0.28); \draw[bond] (Bc.west)--(-0.4,-0.28);
            \draw[bond] (-0.4,0.28) to[out=180,in=180] (-0.4,-0.28);
            % '=' at 1.05 (at 1.15 it nearly touched the left-bulging identity arc)
            \node at (1.05,0) {$=$};
            \draw[bond] (1.5,0.28)--(2.0,0.28); \draw[bond] (1.5,-0.28)--(2.0,-0.28);
            \draw[bond] (1.5,0.28) to[out=180,in=180] (1.5,-0.28);
        \end{scope}
        \node[nlabel] at (1.1,1.35) {left-canonical (trace-preserving)};
        \node[note, text width=3.6cm] at (1.1,0.5) {filled $=A$, open $=\bar A$, arc $=$ identity};

        \draw[draw=black!35, line width=0.4pt] (3.9,-0.1) -- (3.9,5.6);

        % Resolved (\SL{cut panel b}): the old panel (b), the CF/biCF block
        % canonical forms, was removed as fully covered by the cascade figure
        % after the notation sync (same squares, shades, and direct sums
        % there); the caption now points to \cref{fig:cascade} instead.

        %% ===== Panel (b): two pairings (TN loops) =====
        \node[nlabel, text width=3.8cm] at (6.0,5.4) {(b) two pairings on $M_D(\C)$};

        \begin{scope}[shift={(6.0,3.95)}]
            \node[opf] (M) at (-0.65,0) {$X$}; \node[opf] (N) at (0.65,0) {$Y$};
            \draw[bond] (M.east) -- (N.west);
            \draw[bond] (N.east) to[out=50,in=130,looseness=1.6] (M.west);
        \end{scope}
        \node[align=center, text width=3.8cm] at (6.0,3.0)
            {$\mathrm{Tr}(XY)$: \textbf{bilinear}\\ (algebraic single-block FT)};

        \begin{scope}[shift={(6.0,1.65)}]
            \node[opc] (X) at (-0.65,0) {$\bar X$}; \node[opf] (Y) at (0.65,0) {$Y$};
            \draw[bond] (X.north) to[out=90,in=90] (Y.north);
            \draw[bond] (X.south) to[out=-90,in=-90] (Y.south);
        \end{scope}
        \node[align=center, text width=4.0cm] at (6.0,0.6)
            {$\langle X,Y\rangle=\mathrm{Tr}(Y X^\dagger)$: \textbf{sesquilinear}\\ (blocking / spectral; open $=\bar X$)};

    \end{tikzpicture}
    \caption{Two notions that the terms ``canonical form'' and ``inner
        product'' overload in this development: (a) The gauge of the transfer
        operator $E_A(X)=\sum_i A^i X {(A^i)}^\dagger$, as tensor-network
        identities (filled $=A$, open $=\bar A$, an arc $=$ the identity):
        closing the right bonds gives the identity on the left, the unital
        right-canonical gauge $\sum_i A^i {(A^i)}^\dagger=I$; closing the left bonds
        gives the identity on the right, the trace-preserving left-canonical
        gauge $\sum_i {(A^i)}^\dagger A^i=I$. A gauge realizes one or the other;
        both at once, the doubly stochastic gauge discussed in the main text, holds only for a non-generic family.
        (b) The two
        pairings on bond operators the proofs use: the bilinear trace pairing
        $\mathrm{Tr}(XY)$, on which the algebraic single-block argument rests, and
        the sesquilinear Hilbert--Schmidt inner product
        $\langle X,Y\rangle=\mathrm{Tr}(Y X^\dagger)$, which defines the
        transfer-operator adjoint used in the blocking and spectral arguments.}\label{fig:gauge_innerproduct}
\end{figure*}

